\chapter*{ABSTRACT}

Perovskite solar cells are difficult to diagnose from terminal electrical
measurements because mobile ions, recombination, transport-layer interfaces,
contact resistance, and leakage can produce similar changes in a measured
J--V curve. The modeling problem is therefore not only to fit a decreasing
power-conversion-efficiency trajectory. It is to build a device model whose
geometry, parameters, voltage protocol, internal state variables, and outputs
can be inspected and challenged against physical expectations.

This thesis develops a COMSOL-based degradation model for perovskite solar
cells. The electronic device physics are represented with semiconductor
drift-diffusion and Poisson electrostatics in a one-dimensional reduction of
the laboratory p-i-n architecture. From the illuminated side, the stack is
\begin{equation*}
\mathrm{glass}/\mathrm{ITO}/\mathrm{NiO}_x/
\mathrm{MeO\mbox{-}2PACz}/\mathrm{MHP}/
\mathrm{C}_{60}/\mathrm{BCP}/\mathrm{Ag}/
\mathrm{UV\ resin}/\mathrm{glass}.
\end{equation*}
The electrical COMSOL domain
combines NiO$_x$ and MeO-2PACz into a 35 nm hole-transport region, combines
C$_{60}$ and BCP into a 29 nm electron-transport region, represents the
perovskite as a 450 nm absorber, and applies ITO and Ag as electrical boundary
contacts. Encapsulation layers are retained in the physical architecture but
are outside the present electronic transport domain.
Mobile-ion redistribution is implemented through a user-defined drift-diffusion
PDE, and degradation is represented through bounded state variables that map to
COMSOL-accessible parameters such as trap density, ionic charge, series
resistance, shunt resistance, surface recombination, and photogeneration.
Depth-dependent Beer--Lambert generation and preconditioned mobile-ion
profiles are included so that optical absorption and bias-history effects enter
the model as inspectable physical fields rather than empirical terminal
corrections.

The thesis separates slow aging from fast J--V measurement by distinguishing
the stress voltage $V_{\mathrm{stress}}(t)$ from the scan voltage
$V_{\mathrm{app}}(t)$. This separation supports accelerated aging studies,
multi-rate hysteresis simulations, DIT pulse studies for mobile-ion
concentration extraction, and mechanism-isolation studies. The COMSOL
degradation library contains all-off, full coupled, one-mechanism, and
leave-one-out cases so that the J--V signature of each degradation pathway can
be interpreted directly.
Additional architecture-scoped exports add a one-sun temperature-aging J--V
sweep and a 6 $\times$ 6 light-temperature internal-state profile grid. These
outputs connect terminal photovoltaic metrics to spatial summaries of
$u_{\mathrm{ion}}$, $D_{\mathrm{trap}}$, resistance damage, effective trap
density, recombination, and electric field across controlled stress
conditions.

The experimental comparison is qualitative. The COMSOL baseline is more ideal
than the experimental X12 baseline, while the accelerated light-plus-heat case
reproduces the direction of degradation, including PCE reduction, fill-factor
loss, and J--V knee collapse. The current implementation demonstrates
qualitative mechanism separation and produces surrogate-ready COMSOL state
summaries, but experimental calibration remains ongoing. Future work should
improve absolute calibration, validate ion-sensitive parameters, implement
light-soaking and self-healing terms, complete corrected low-light J--V
exports, and develop architecture-specific parameter sets.

\clearpage
